% Abstract (v2, trade-off framing). Every claim is recoverable from Tables tab:decon, tab:gains,
% tab:disclosure-units, tab:identity or Section 6. At most 200 words (ACL formatting guidelines):
% count the rendered text after any change (scripts/papers/check_paper_limits.py).
\begin{abstract}
A character wiki that reveals only what a reader has reached invites decomposing the novel
into a graph of claims filtered by volume. We instrument this design as \sys and test it
against six systems at every cutoff of a five-volume series
(\res{T1-anne-gold-t5} gold facts; LLM-finalized labels\ifnwhuman{}; two human reviewers
relabeled \res{HREV-total} original-text items\fi{}). A renamed and paraphrased copy (the
treated text) tests reliance on recognition: closed-book recall falls
from \respt{M1-anne-avg-X1} to \respt{M1-annedecon-avg-X1}, while systems given text lose
at most \res{M1-annedecon-maxloss}. The design trades context for accumulation. \sys's
pages grow with the story: on treated text it best covers the most-documented characters
(recall \respt{CEN-annedecon-B5}; long context \respt{CEN-annedecon-B1}; original:
\respt{CEN-anne-B5}, \respt{CEN-anne-B1}), loses one fact across cutoffs per text (others:
\res{DROP-anne-B2}--\res{DROP-annedecon-B2}), and on the original text publishes 3.8 times
long context's distinct assertions. But claims lose context: \sys turns a prediction into a
marriage (treated text), groups entities unlike prefix-built indexes
(\respt{ID-anne-strict-t1} of first-cutoff pages, original text) and misplaces true
statements. Prompting with only the text read discloses no assessed later fact; an
instruction to withhold later volumes discloses \res{M4U-annedecon-X2} later facts on
treated text (original: \res{M4U-anne-X2}). Code, gold facts and labels: \anonrepo.
\end{abstract}
